\section{Introduction}\label{sec:introduction}

Trajectory data represents object movement as sequences of timestamped locations and supports mobility analysis, route retrieval, and regional activity studies.
  Key-value (KV) stores such as HBase provide scalable storage through RowKey lookups and ordered range scans.
  A trajectory query system must map spatial, temporal, and attribute predicates to these physical access mechanisms, then refine retrieved candidates to preserve the query semantics.
  Its performance and correctness therefore depend on data layout, index organization, and access-plan selection.

Large language models (LLMs) make it possible to express database requests in natural language.
xx et al. have studied text-to-SQL and text-to-NoSQL approaches, and there are also numerous related works on LLM-based query optimization, such as: TBD, TBD
However, these systems and methods have not yet addressed physical-layer issues related to trace queries on key-value (KV) stores, including RowKey ranges, dedicated trace indexes, and block-level candidate operations. This gives rise to the following three challenges.

\noindent\textbf{Natural-language translation.} Requests may omit or ambiguously express temporal, spatial, attribute, or similarity constraints.
  A system must resolve intent against supported trajectory semantics and its logical catalog, and request clarification when required information is missing.

\noindent\textbf{Access-path selection.} A query may use or combine temporal, spatial, and attribute indexes, followed by candidate operations and exact filtering.
  The best path depends on both query selectivity and the costs of scans, retrieval, and refinement; fixed rules may not adapt to these differences.

\noindent\textbf{Plan correctness.} A syntactically valid plan can still omit a relevant index range, change query semantics, or compile to an invalid HBase request.
  Exact filtering cannot recover records missed during access, so plans require validation for semantics, candidate coverage, and physical safety.


\begin{figure}[tbp]
\centering
\resultplaceholder{TBD: Diagrams should be provided for the three challenges.}{
Specifically, there are three subdiagrams—a, b, and c—which should have similar structures and correspond to the three challenges, respectively. Challenges 2 and 3 occur during the query phase, while Challenge 1 occurs during the processing phase.}
\caption{TBD}
\label{fig:kart-motivation}
\end{figure}

To address these limitations, we present \KART{}, a natural-language trajectory query system for KV stores.
  It maps requests through schema-constrained DraftIR and deterministically bound BoundIR representations, searches a bounded space of legal access plans using cost estimates, and validates compiled candidates before execution.
  Two offline LLM-assisted paths curate scoped rewrite rules and calibrate plan-cost estimates; runtime proposals remain constrained by the legal operator space, and the cost-based path remains available when a rewrite is rejected or unavailable.
  In summary, we make the following contributions:

\begin{itemize}
\item We design an end-to-end query framework that separates language interpretation from binding and physical access planning through typed DraftIR and BoundIR representations.
\item We develop bounded, cost-guided access-path search with two offline LLM-assisted artifacts: conditional rewrite rules for candidate generation and cost patches for plan ranking.
\item We have built a multi-layered correctness assurance mechanism. The verification process covers the entire workflow, from semantics to concrete plans. Only plans that pass verification are compiled into HBase operations.
\item We conducted extensive experiments on two real-world datasets of different sizes, demonstrating KART's query efficiency and robustness across various scenarios and levels of complexity.
\end{itemize}
